% ============================================================
% Chapter 6: Experimental Design
%   Section: Construction of the Custom Validation Set
%     Subsection: Annotation Procedure
% Included from chapters/06-experimental-design/03-custom-validation-set/custom-validation-set.tex
% ============================================================

\subsection{Annotation Procedure}
\label{subsec:custom-procedure}

Claims are taken from articles published by the positive-news and
science sources the system ingests, and neither the articles nor the
claims are chosen by the annotator. Each day a batch of about ten
articles is drawn at random from the sources' current feeds and topic
pages: at most one article per source, the two languages alternated,
every article already labelled or previously proposed excluded, and
articles whose predicted topic belongs to an under-filled topic group
preferred, after articles published within the previous sixty days,
which keep small the gap between the evidence admissible under
Rule~3 and the present-day web the system searches. The draw is seeded
with the date, so the same set of candidate links yields the same
draw.
For each article, the system's own claim extractor and anchor-claim
selector propose between one and three claims, which are therefore
the claims the system would itself verify.

The annotator either labels a proposed claim or rejects it with a
recorded reason: an opinion or interpretation, a prediction, a
checkable but trivial statement, a fragment of a sentence, or a
duplicate. Those rejections are kept, and the share of proposals
accepted is reported as the precision of the system's claim
selection as judged by the annotator. Each record holds one atomic
claim, kept in the article's own wording; a proposed sentence that
asserts two things is split into two records, following the claim
decomposition practice of real-world fact-checking datasets
\cite{guo2022survey}. The annotator then searches for evidence
independently of the system, records the links on which the verdict
rests, and assigns the label under the rules of
Section~\ref{subsec:custom-rules}. The system's verdict is never
shown during annotation.

Records are entered through a small web form built for this purpose
and kept apart from the system itself: a single program using only the
Python standard library, so that labelling depends neither on the
system's services nor on any installed package. Its fields are exactly
the keys of Table~\ref{tab:custom-format}, and it enforces the rules
that can be checked mechanically at the moment a record is saved:
a verdict other than \texttt{UNVERIFIED} requires at least one
evidence link; a claim backed only by its protagonist's own source
must be labelled \texttt{UNVERIFIED} and the source named; and an
evidence date later than the article's publication date is rejected.
The same form shows the running count per verdict and topic group
against the targets of Section~\ref{subsec:custom-size}, which is how
balance is steered during annotation rather than assessed after it.

Each record is stored as its own file (\texttt{fact001.json},
\texttt{fact002.json}, \ldots) in the project's version-controlled
repository. Every verified claim is therefore individually visible,
and every later correction is recorded as a change to that one file,
which leaves an audit trail of the annotation. Once the set is
complete, the files are validated again and merged, in order, into a
single JSON Lines file in the X-Fact format, which is what the
evaluation harness reads. The per-record files remain the source, and
the merged file is derived from them.
